\chapter{Journaux et reproduction des figures}\label{ann:journaux}

Ce qui produit un chiffre publié est commité : outils et journaux. Cette annexe dit où sont les journaux et quelle donnée lit chaque graphe. Les deux tables sont engendrées par \fichier{memoire/outils/annexes.py}.

\section{Reconstruire le mémoire}

Depuis \fichier{memoire/}, dans Git Bash, avec MiKTeX (lualatex, biber, latexmk) et Python 3.13 sans dépendance :
\begin{lstlisting}[language=bash, numbers=none]
./construire.sh          # données des figures, annexes, puis les deux PDF
./construire.sh fr       # une seule langue
./construire.sh propre   # efface les fichiers intermédiaires
\end{lstlisting}
Chaque graphe pgfplots lit un CSV commité ; chaque schéma TikZ est un fichier de \fichier{memoire/figures/schemas/} ; les commandes des figures sont communes aux deux langues, qui ne diffèrent que par leurs textes (\verb|\EF{anglais}{français}|). Le texte d'un mémoire se lit en parallèle de l'autre : mêmes chapitres, mêmes étiquettes, mêmes figures.

\section{Les journaux}

\input{\Racine/annexes/journaux}

\section{Les données de chaque graphe}

\input{\Racine/annexes/figures}
